% paper/sections/07_advection.tex
%
% §6.1  CLS $\psi$ 用 FCCD と運動量用 UCCD6 (per-term spatial discretization)
%       §6 chapter intro は 07_0_scheme_per_variable.tex に移行済み (CHK-223).
%       本ファイルは §6.1 subsection から始まる．
%
% backward-compat: §6 chapter intro 時代の section ラベルを保持
\phantomsection\label{sec:discretization_modules}% legacy chapter alias
\phantomsection\label{sec:advection}% legacy chapter alias

% ─────────────────────────────────────────────────────────────────────────────
\subsection{移流項：CLS $\psi$ 用 FCCD と運動量用 UCCD6}
\label{sec:advection_motivation}
% ─────────────────────────────────────────────────────────────────────────────

\noindent\textbf{Part 1 連続定式化との関係：}
本節は §~\ref{sec:cls_advection}（§3.2 CLS 保存形移流方程式，
式~\eqref{eq:cls_advection}）と §~\ref{sec:onefluid}（§2 One-Fluid NS，
式~\eqref{eq:NS_onefluid}）の運動量移流項 $(\bm{u}\cdot\bnabla)\bm{u}$ の
空間離散化を統一して扱う．
CLS 移流方程式の再掲：
\begin{equation}
  \frac{\partial\psi}{\partial t} + \bnabla\cdot(\psi\bm{u}) = 0
  \label{eq:cls_advection_ref}
\end{equation}
ch13 production stack では，CLS $\psi$ には面値保存型の FCCD
（§~\ref{sec:fccd_def}）を，運動量 $(u,v)$ には演算子内部散逸の
UCCD6（§~\ref{sec:uccd6_def}，L1/L2/L3 既定）を用いる．
FCCD による BF 整合移流 Option B/C は §~\ref{sec:fccd_advection}（§6.2）で展開し，
界面張力項は §~\ref{sec:fccd_bf_sub_system}（§8）で扱う．

% ─────────────────────────────────────────────────────────────────────────────
\subsubsection{移流スキーム選択の設計根拠：WENO5 不適合と CLS に FCCD を採る理由}
\label{sec:advection_weno5_critique}
% ─────────────────────────────────────────────────────────────────────────────

WENO5~\cite{JiangShu1996} は圧縮性流体力学の衝撃波捕捉向けに設計されたスキームであり，
CLS 界面移流には以下の 3 点で非適合である
（詳細な定式化と性質は付録~\ref{app:weno5} を参照）：

\begin{enumerate}
  \item \textbf{過剰散逸：}
        CLS の $\psi = H_\varepsilon(\phi)$ は滑らかな関数であり真の不連続を含まないが，
        WENO5 の滑らかさ指標はこの勾配を「不連続の兆候」と誤検知し，
        不必要な非線形散逸を常時付加する．
        結果として界面の微細構造（液膜・せん断巻き込み）が非物理的に平滑化される．

  \item \textbf{Balanced-Force 精度の侵食：}
        WENO5 の非線形重みが $\psi$ プロファイルに振動・ぼやけを残すことで，
        Balanced-Force 条件（§~\ref{sec:balanced_force}）が確保する
        寄生流れ $\Ord{h^6}$ の高精度を実用上無効化する危険がある．

  \item \textbf{ステンシル非整合：}
        本フレームワークは全空間離散化に CCD の 3 点コンパクトステンシルを用いるが，
        WENO5 の 5 点ステンシルを混在させると
        打ち切り誤差の非整合（$\Ord{h^5}$ が $\Ord{h^6}$ を律速）が生じ，
        境界条件処理の一元化も失われる．
\end{enumerate}

\noindent
\textbf{結論}：純 WENO5 単独は CLS 界面移流に不適合であるため，
本稿は CLS 移流に \textbf{FCCD（Flux-form CCD；面値保存型，§~\ref{sec:fccd_def}）}を採用する．
すなわちセル境界 $x_{i+1/2}$ における面値フラックス $(\psi u)_{f,i+1/2}$ を
Hermite 補間で構成し，その差分でセル中心の発散
$F_i = [(\psi u)_{f,i+1/2} - (\psi u)_{f,i-1/2}]/h$ を求める
（§\ref{sec:cls_stages} Stage A，ch13 production の FCCD 面ジェット移流）．
FCCD は\textbf{保存型 FVM のテレスコーピング構造}を持ち，
$\psi$ プロファイル kink での節点 CCD overshoot を経由しないため，
界面の急峻な遷移層に対しても安定かつ厳密質量保存的である．
WENO5 は \emph{比較参照スキーム}として付録~\ref{app:weno5} に収録するに留め，
本稿の主実装には用いない（§\ref{sec:verify_weno5_vs_dccd} で FCCD と WENO5 を定量比較）．\footnote{ch13
production stack 確立前の旧実装では帯状散逸付き節点 CCD を CLS 移流に使用していた．
当該手法は線形フィルタによる近似的安定化に依存し，非周期境界での厳密質量保存を満たさない．
FCCD への移行は保存型構造による厳密質量保存と界面安定性の両立を理由とする．
旧実装の理論的詳細と歴史的経緯は §~\ref{sec:dissipative_ccd}（§4.4）を参照．}

% ─────────────────────────────────────────────────────────────────────────────
\subsubsection{気液界面移流の困難と純 CCD の限界}
\label{sec:advection_ccd_limit}

$\psi$ は界面付近で $0$ から $1$ へと急激に変化する（界面幅 $\sim 2\varepsilon$）．
この急峻なプロファイルを移流方程式~\eqref{eq:cls_advection_ref}で輸送するとき，
\textbf{純中心差分を前進 Euler 型の時間積分と組み合わせると線形安定性を失う}．
虚軸安定域を持つ Runge--Kutta 法では許容 CFL が存在しうるため，
以下ではまず最も単純な前進 Euler との組合せで純 CCD の限界を確認する．

フォン・ノイマン安定性解析による確認：
CCD は純中心差分（修正波数 $k^*(kh)$ が実数）であるため，
移流方程式 $\partial\psi/\partial t + u\,\partial\psi/\partial x = 0$ に
前進 Euler 時間積分を組み合わせると，Fourier モード $\psi_i = e^{i\xi i}$ に対して
増幅係数は
\begin{equation}
  |g(\xi)|^2 = 1 + \sigma^2 \bigl[\hat{k}^*(\xi)\bigr]^2 > 1
  \quad (0<\xi<\pi,\ \hat{k}^*(\xi)\neq0), \;\sigma = \frac{u\Delta t}{h} > 0
  \label{eq:ccd_adv_instability}
\end{equation}
ここで $\xi = kh \in (0,\pi]$ は無次元波数，
$\hat{k}^*(\xi) = k^*(\xi)\,h$ は CCD の無次元修正波数（実数；§~\ref{sec:CCD} の Fourier 解析参照），
$\sigma$ は CFL 数である．
ナイキスト点 $\xi=\pi$ では中心型微分がそのモードを感知できず中立となるが，
それ以外の可視波数では $|g|>1$ となる．
したがって純 CCD + 前進 Euler の組合せは，CFL を小さくしても増幅率を1未満にはできない．
TVD-RK3 を使う場合の線形安定性は，後述の虚軸安定域条件で別途判定する．

この不安定性の物理的帰結は \textbf{Gibbs 振動}（$\psi < 0$ や $\psi > 1$
の非物理的値の発生）であり，界面のにじみおよび密度場の不整合として現れる
（第~\ref{sec:intro}章~§~\ref{sec:failure_modes}）．

\subsubsection{FCCD 面値演算子の採用：保存型フラックスによる安定化}
\label{sec:advection_dccd_design}% backward-compat: 旧 DCCD 設計 \ref ラベル維持

第~\ref{sec:fccd_def}節で定式化した FCCD（Flux-form CCD）面値演算子
$\mathbf{P}_f$ を CLS 移流に適用する．
基本アイデアは，節点 CCD で $\partial(\psi u)/\partial x|_i$ を直接評価する代わりに，
\textbf{セル境界 $x_{i+1/2}$ における面値フラックス $(\psi u)_{f,i+1/2}$ を
Hermite 補間 $\mathbf{P}_f$ で構成し，その差分でセル中心の発散を求める}ことである．

\medskip
\noindent\textbf{移流への適用．}
フラックス場 $(\psi u)_i = \psi_i u_i$ とその CCD 微分 $\partial_x(\psi u)_i$ を
入力として面値 $(\psi u)_{f,i+1/2}$ を構成し，セル中心発散をとる：
\begin{align}
  (\psi u)_{f,i+1/2} &= \mathbf{P}_f\!\left[\,(\psi u)_i,\, \partial_x(\psi u)_i,\, (\psi u)_{i+1},\, \partial_x(\psi u)_{i+1}\,\right]
  \quad\text{（FCCD 面値，$\Ord{h^6}$）}
  \label{eq:dccd_adv_ccd} \\[4pt]
  F_i &= \frac{(\psi u)_{f,i+1/2} - (\psi u)_{f,i-1/2}}{h}
  \quad\text{（保存型発散）}
  \label{eq:dccd_adv_filter}
\end{align}
得られた発散 $F_i$ を時間積分（§~\ref{sec:time_int}）の
空間演算子として使用する：$\mathcal{L}(\psi) = -F$
（PDE $\partial\psi/\partial t = -\bnabla\cdot(\psi\mathbf{u})$ の符号に整合）．
ch13 production の FCCD 面ジェット移流はこの定式化を実装する．

\medskip
\noindent\textbf{安定化機構：面値評価が振動源を切る．}
節点 CCD では $\psi$ プロファイルの kink を 3 点コンパクトステンシルが
直接横切るため，急峻な遷移層で $\partial\psi/\partial x|_i$ の補間が overshoot し，
それが §~\ref{sec:advection_ccd_limit} で示した Gibbs 振動を励起する．
FCCD は面値 $(\psi u)_{f,i+1/2}$ を Hermite 補間で滑らかに構成し，
セル中心の発散をその差分（式~\eqref{eq:dccd_adv_filter}）で求めるため，
振動源となる節点での直接微分を経由しない．
これが「面値保存型」アプローチの安定化機構であり，
線形フィルタによる高波数応答低減（近似的安定化；§~\ref{sec:dissipative_ccd}）とは構造的に異なる．
詳細な精度・分散関係は §~\ref{sec:fccd_def} 参照．

\medskip
\noindent\textbf{質量保存性．}
FCCD は\textbf{保存型 FVM のテレスコーピング構造}を持つ：
セル境界面値 $(\psi u)_{f,i+1/2}$ は隣接セル $i,\,i+1$ で共有されるため，
領域全体での $F_i\,h$ の総和は境界フラックスのみで決まる：
\[
  \sum_i F_i\,h = (\psi u)_{f,N+1/2} - (\psi u)_{f,1/2}
\]
周期境界では右辺がゼロとなり厳密に質量保存．
非周期境界・非一様格子・界面適合格子でもテレスコーピング構造が保たれるため，
従来の節点 CCD 微分に伴う「近似的保存」（周期境界以外で破れる近似的成立）が
FCCD では構造的に解消される．
これが ch13 production で FCCD を採用した第一の理由である
（WENO5 保存形との比較：付録~\ref{app:weno5}；
定式化§~\ref{app:weno5_formulation}，再構成係数§~\ref{app:weno5_coef}，Lax-Friedrichs 分割§~\ref{app:lf_splitting}）．

\medskip
\noindent\textbf{TVD 性と clamp 操作．}
FCCD は線形 Hermite 補間に基づくため一般には TVD ではなく，
急峻な界面では $\psi \notin [0,1]$ の overshoot/undershoot が小幅に発生しうる．
本稿では各 RK3 ステージ完了後に値域制限を課す：
\begin{equation}
  \psi_i \leftarrow \max(0,\,\min(1,\,\psi_i))
  \label{eq:psi_clamp}
\end{equation}
\label{warn:adv_clamp}%
これは $\psi \in [0,1]$ の物理的制約を強制する非線形操作であり，
一般には質量保存形ではない．実運用では clamp 後に
§~\ref{sec:dccd_mass_conservation} の界面重み付き質量補正を適用し，
値域制限で生じた総量誤差を戻す
（完全なアルゴリズムは §~\ref{sec:algo_flow}, Step 1）．

\begin{tcolorbox}[warnbox, title={線形 CFL 上界の目安（FCCD + TVD-RK3，線形理論・未検証）}]
\label{result:dccd_stability}% backward-compat: 旧 DCCD 安定性 ラベル維持
線形定係数移流に対するフォン・ノイマン解析では，
TVD-RK3 の虚軸安定域 $|z|\le\sqrt{3}$ と
FCCD 面値演算子の最大修正波数
$\max_\xi |\hat{k}^*_{\mathrm{FCCD}}(\xi)| \approx 2.0$
から，CFL 上界の目安は $\sigma_{\max} \approx \sqrt{3}/2.0 \approx 0.87$ となる．
本稿の運用 CFL $\leq 0.5$（式~\eqref{eq:dt_adv}）はこの上界を下回る．
ただしこれは定係数・周期境界・局所凍結係数近似での線形評価であり，
非線形 CLS 移流・値域制限・非一様格子・境界条件を含む実計算の安定性は
§~\ref{sec:verify_cls_advection} の数値検証で確認する．
\end{tcolorbox}

\label{alg:dccd_adv}%
完全な CLS 移流の離散アルゴリズム（TVD-RK3 + FCCD 面値演算子 + 値域制限 + 質量補正）は
§~\ref{sec:algo_flow}（Step 1）にまとめる．

\paragraph{誤差指標の注意：$\psi$ 空間 vs $\phi$ 空間}
CLS 移流の形状誤差を $L_2(\psi - \psi_0)$ で評価する場合，
$\psi$ の遷移層幅 $\varepsilon$ に依存するバイアスが生じる：
界面厚 $\varepsilon$ が大きいほど $\psi$ 差が空間的に広がり，
$L_2(\psi)$ が見かけ上悪化する．
異なる $\varepsilon/h$ 設定間の比較や，
格子収束テストの物理的解釈には，
ロジット逆変換 $\phi = \varepsilon\ln(\psi/(1-\psi))$
で $\phi$ 空間に射影した $L_2(\phi - \phi_0)$ を用いるべきである．
$\phi$ 空間誤差は界面位置の幾何学的ずれを直接反映し，
$\varepsilon$ に依存しない．
本稿の実験結果（§~\ref{sec:verify_cls_advection}）では
$L_2(\psi)$ と $L_2(\phi)$ の両方を報告する．

\begin{tcolorbox}[mybox, title={本稿の空間離散化スキームの役割分担}]
\label{box:scheme_roles}
\begin{center}
\renewcommand{\arraystretch}{1.3}
\begin{tabular}{>{\raggedright\arraybackslash}p{2.5cm}>{\raggedright\arraybackslash}p{3.8cm}>{\raggedright\arraybackslash}p{1.5cm}>{\raggedright\arraybackslash}p{4cm}}
\hline
対象 & スキーム & 精度 & 根拠 \\
\hline
CLS 移流（$\psi$） & FCCD（面値保存型；§~\ref{sec:fccd_def}，§~\ref{sec:advection_motivation}，§~\ref{sec:cls_stages} Stage A） & 面値 $\Ord{h^6}$；保存型 FVM テレスコーピング & 面値評価で節点 CCD overshoot を回避；非周期境界・非一様格子でも厳密質量保存；WENO5 は比較参照（付録~\ref{app:weno5}） \\
運動量移流（$u,v$） & UCCD6（§~\ref{sec:uccd6_def}） & CCD 微分 $\Ord{h^6}$；ハイパー粘性 $\Ord{\sigma h^7}$（次数保存） & 演算子内部散逸；$\omega_2^8$ 選択的 Nyquist（L2/L3 bulk 既定） \\
CLS 再初期化 & Ridge-Eikonal（§~\ref{sec:ridge_eikonal}，§~\ref{sec:cls_compression}） & Eikonal 距離 $\Ord{h^6}$ + Ridge ガウス射影 & 反復不要の 1 ショット再構成；§3.4 Part 1 連続定式化の離散実装 \\
曲率・速度微分 & CCD（§~\ref{sec:CCD}） & $\Ord{h^6}$ & 高精度・コンパクト \\
速度・圧力場ノイズ除去 & 適応制御散逸演算子（§~\ref{sec:dissipative_ccd}） & — & $S(\psi)$ で界面保護 \\
圧力 Poisson & 欠陥補正法 PPE + CCD 勾配（§~\ref{sec:pressure}） & $\Ord{h^5}$--$\Ord{h^7}$† & Balanced-Force 整合性（DC k=3） \\
\hline
\end{tabular}
\end{center}
{\small †\ DC k=3 では Dirichlet $\Ord{h^7}$，Neumann BC $\Ord{h^5}$（§~\ref{sec:verify_split_ppe_dc}）．速度補正の圧力勾配は CCD $\Ord{h^6}$（§~\ref{sec:pressure_accuracy_summary}）．}
\end{tcolorbox}
